% !TeX root = ../Cosmology.tex


%% --------------------------------------------------------
\chapter{Інфлатон і його потенціал}
\label{sec:inflaton}
%% --------------------------------------------------------

У попередньому розділі ми встановили, що інфляція~---
епоха прискореного розширення з $N \gtrsim 60$ $e$-folds~---
\emph{кінематично} вирішує проблеми горизонту і плоскості.
Але залишилось відкритим головне питання: \emph{що} фізично
рухає це розширення?

З другого рівняння Фрідмана~\eqref{eq:F2} відомо, що умова
$\ddot{a} > 0$ вимагає речовини з рівнянням стану $w < -1/3$
(§\,\ref{sec:kinematics_bridge}). Жодна відома форма матерії~---
ні газ, ні випромінювання, ні темна матерія~--- цього не
задовольняє: всі вони мають $w \geq 0$.

Вихід пропонує квантова теорія поля. У сучасній фізиці
фундаментальні взаємодії описуються не частинками, а
\emph{полями}~--- величинами, що задані в кожній точці
простору-часу. Електромагнітне поле, поле Хіггса, кваркові
поля~--- це все приклади. Найпростіший тип~--- \emph{скалярне
поле} $\phi(\mathbf{x},t)$: в кожній точці одне число, без
напрямку (на відміну від векторного електричного поля).

Виявляється, скалярне поле з певним потенціалом $V(\phi)$
здатне мати $w \approx -1$~--- якщо воно повільно еволюціонує
вздовж пологого схилу $V(\phi)$. Таке гіпотетичне поле
називають \textbf{інфлатоном}. Його динаміці і присвячений
цей розділ.

%\medskip
%\noindent\textbf{Логічна карта розділу:}
%$V(\phi) \;\longrightarrow\; \varepsilon,\,\eta_V \;\longrightarrow\;
%a(t) \sim e^{Ht} \;\longrightarrow\; \mathcal{P}_{\mathcal{R}}(k,n_s,r)
%\;\longrightarrow\; \text{розігрів}$.

\medskip
\noindent\textbf{Логічна карта розділу}

\begin{center}
\begin{tikzpicture}[roadmap]

\node[box] (V) {$V(\phi)$};
\node[box, right=of V] (params) {$\varepsilon,\eta_V$};
\node[box, right=of params] (expansion) {$a\sim e^{Ht}$};
\node[box, right=of expansion] (models) {Моделі};
\node[box, right=of models] (spectrum) {$n_s,r$};
\node[box, right=of spectrum, fill=red!5, draw=red!70!black] (reheat) {Розігрів};

\draw[arrow] (V) -- (params);
\draw[arrow] (params) -- (expansion);
\draw[arrow] (expansion) -- (models);
\draw[arrow] (models) -- (spectrum);
\draw[arrow] (spectrum) -- (reheat);

% Підписи зверху
\node[font=\tiny, gray, above=0.1cm of V] {Двигун};
\node[font=\tiny, gray, above=0.1cm of params] {Slow-roll};
\node[font=\tiny, gray, above=0.1cm of expansion] {Розширення};
\node[font=\tiny, gray, above=0.1cm of models] {Зоопарк};
\node[font=\tiny, gray, above=0.1cm of spectrum] {Спостереження};
\node[font=\tiny, gray, above=0.1cm of reheat] {Естафета};

\end{tikzpicture}
\end{center}

\medskip
\noindent\textbf{Крок 1: $V(\phi)$~--- двигун інфляції.}
Скалярне поле інфлатона $\phi$ із потенціалом $V(\phi)$~---
найпростіший механізм, здатний забезпечити від'ємний тиск
$p \approx -\rho$. Ключова умова: поле має повільно «котитися»
по пологому схилу, щоб кінетична енергія залишалася малою
($\dot{\phi}^2 \ll V$). Форма потенціалу~--- головна
невідомість теорії.

\medskip
\noindent\textbf{Крок 2: $\varepsilon, \eta_V$~--- параметри повільного кочення.}
Два безрозмірних числа повністю контролюють динаміку:
$\varepsilon \equiv -\dot{H}/H^2 < 1$ гарантує прискорення,
$|\eta_V| \ll 1$ гарантує тривалість. Обидва виражаються
через похідні потенціалу~--- тому форма $V(\phi)$ прямо
визначає спостережувані наслідки.

\medskip
\noindent\textbf{Крок 3: $a(t) \sim e^{Ht}$~--- розв'язання проблем.}
У режимі slow-roll $H \approx \mathrm{const}$ і простір
розширюється експоненційно. За $N \sim 50$--$60$ $e$-folds
кривина розгладжується, горизонтна проблема знімається.
Але головне~--- народжуються квантові флуктуації.

\medskip
\noindent\textbf{Крок 4: зоопарк моделей.}
Існує понад 100 конкурентних моделей інфляції. Ми розглянемо
три ключові: хаотичну інфляцію Лінде ($V\propto\phi^2$,
виключена), модель Старобінського ($R^2$-гравітація, «золотий
стандарт») і моделі зі зламаним спектром (мотивовані JWST).

\medskip
\noindent\textbf{Крок 5: $\mathcal{P}_{\mathcal{R}}(k)$~--- міст до спостережень.}
Форма потенціалу на вузькій ділянці ($\sim 7$--$8$ $e$-folds
із $50$--$60$) кодується у первинному спектрі потужності
через два числа: спектральний індекс $n_s$ і
тензорно-скалярне відношення $r$. Саме їх вимірює Planck~---
і саме вони відсівають або підтверджують моделі.

\medskip
\noindent\textbf{Крок 6: розігрів~--- передача естафети.}
Коли поле досягає мінімуму $V(\phi)$, воно починає
осцилювати і розпадатися на частинки Стандартної моделі.
Квантовий вакуум перетворюється на гарячу плазму~---
і починається термічна історія Всесвіту, описана в наступних
розділах.

%\bigskip
%З другого рівняння Фрідмана~\eqref{eq:F2} умова прискореного
%розширення $\ddot{a} > 0$ вимагає $w < -1/3$
%(§\,\ref{sec:kinematics_bridge}).
%Жодна відома форма звичайної матерії цього не задовольняє~---
%тому вводять інфлатон.

%% --------------------------------------------------------
\section{Механізм та рівняння динаміки інфлатона}
\label{subsec:inflaton_dynamics}
%% --------------------------------------------------------


\begin{keybox}{Як народилася ідея інфлатона: від фазових переходів до slow-roll}

Алан Гут прийшов до інфляції у 1980--1981 роках не через проблему
пласкості як таку, а через фізику \emph{фазових переходів} у теоріях
Великого об'єднання (GUT). При температурі $T \sim 10^{15}$~ГеВ
симетрія GUT розбивається: хіггсівське поле теорії «вибирає» мінімум
потенціалу, як вода замерзає в лід. Якщо цей перехід є переходом
\emph{першого роду}~--- через нуклеацію бульбашок істинного вакууму
в морі хибного~--- то хибний вакуум має від'ємний тиск $p = -\rho$
і ненульову густину енергії, яка не розбавляється розширенням.
Це і є рівняння стану $w = -1$, що забезпечує $\ddot{a} > 0$.
Гут назвав це \emph{репульсивною гравітацією хибного вакууму}~\cite{guth_1981}.

Однак у тій самій статті він відкрито визнав фундаментальну ваду
власної моделі~--- \textbf{проблему виходу} (\emph{graceful exit
problem}): бульбашки істинного вакууму нуклеюють і розширюються
зі швидкістю світла, але простір між ними розширюється швидше~---
експоненційно. Бульбашки ніколи не встигають злитися в однорідний
гарячий Всесвіт. Результат~--- катастрофічно неоднорідна картина,
протилежна до спостережуваної.

Розв'язок прийшов у 1982 році незалежно від Лінде~\cite{linde_1982}
та Олбрехта з Стейнхардтом~\cite{albrecht_1982}: якщо відмовитися
від тунелювання через бар'єр і замінити його \emph{повільним
коченням поля по пологому плато}~--- перехід другого роду~--- то
поле природно досягає мінімуму, осцилює і розігріває Всесвіт.
Проблема виходу знімається автоматично.
Саме ця ідея~--- \emph{slow-roll}~--- і є предметом цього розділу.
\end{keybox}

Найпростішою реалізацією такого механізму в теорії поля є однопольова
модель, де прискорене розширення забезпечується однорідним скалярним
полем \emph{інфлатоном} $\phi(t)$, яке повільно «котиться» вздовж
пологого потенціалу $V(\phi)$ (див. рис.~\ref{fig:slow_roll}). Фізичну аналогію можна провести з
рухом важкої кульки по пологому схилу за наявності сильного
в'язкого тертя: кулька рухається повільно і рівномірно, не
набираючи кінетичної енергії. Завдяки просторовій однорідності
фону, на цьому етапі просторовими градієнтами поля можна знехтувати.

%---------------------------------------------------------
\begin{figure}[h!]
\centering
\begin{tikzpicture}[scale=1.0]
  % Осі координат
  \draw[->, >=stealth, thick] (-0.5, 0) -- (4.5, 0) node[right] {$\phi$};
  \draw[->, >=stealth, thick] (0, -0.5) -- (0, 4.0) node[above] {$V(\phi)$};

  % Крива потенціалу
  % Функція: 0.2*(x-2.5)^3 + 0.5*(x-2.5)^2 + 0.3
  \draw[theoryred, ultra thick, samples=100, domain=0.2:4.2]
    plot (\x, {0.2*(\x-2.5)^3 + 0.6*(\x-2.5)^2 + 0.3});

  % Положення кульки (задаємо координату x, наприклад x = 0.8)
  \def\ballx{0.8}

  % Отримуємо координату для малювання кульки (додаємо +0.18 до Y, щоб підняти її НАД кривою)
  \coordinate (ball) at (\ballx, {0.2*(\ballx-2.5)^3 + 0.6*(\ballx-2.5)^2 + 0.3 + 0.18});

  % Малюємо кульку
  \shade[ball color=accent] (ball) circle (0.18);

  % Позначення мінімуму (вакууму), якщо потрібно
  \draw[dashed, gray] (2.5, 0) -- (2.5, 0.3);
  \node[below, font=\scriptsize] at (2.5, 0) {$\phi_{\min}$};
\end{tikzpicture}

\caption{Схематичне зображення повільного скочування (slow-roll) інфлатонного поля $\phi$ у потенціалі $V(\phi)$ до мінімуму $\phi_{\min}$.}
\label{fig:slow_roll}
\end{figure}
%---------------------------------------------------------


Дія для скалярного поля $\phi$ у просторі FLRW:
\begin{equation}\label{eq:inflaton_action}
  S = \int d^4x\,\sqrt{-g}\left[\frac{1}{2}g^{\mu\nu}\partial_\mu\phi\,\partial_\nu\phi - V(\phi)\right].
\end{equation}
Для однорідного поля $\phi(t)$ маємо $\sqrt{-g} = a^3$ і $g^{00} = -1$,
тому підінтегральний вираз зводиться до:
\begin{equation*}
  \mathcal{L} = a^3\!\left[\frac{\dot{\phi}^2}{2} - V(\phi)\right].
\end{equation*}
Рівняння Ейлера--Лагранжа $\frac{d}{dt}\frac{\partial\mathcal{L}}{\partial\dot{\phi}} - \frac{\partial\mathcal{L}}{\partial\phi} = 0$ дають:
\begin{equation*}
  \frac{d}{dt}\!\left(a^3\dot{\phi}\right) + a^3 V'(\phi) = 0,
\end{equation*}
що після ділення на $a^3$ зводиться до\footnote{Це рівняння є не чим іншим, як коваріантним рівнянням Кляйна-Ґордона-Фока $\left(\Box + m^2\right)\phi = 0$ для вільного поля (або з урахуванням самодії через потенціал $V'(\phi)$), записаним у FRLW-метриці для однорідного скалярного поля $\phi(t)$.}:
\begin{equation}\label{eq:klein_gordon}
    \ddot{\phi} +3H\dot{\phi} + V'(\phi) = 0,
\end{equation}
де $3H\dot{\phi}$ виникає з диференціювання $a^3$:
$\frac{d}{dt}(a^3\dot{\phi}) = a^3(\ddot{\phi} + 3H\dot{\phi})$~---
це і є те саме «хабблівське тертя», яке фізично відіграє роль
в'язкого опору у аналогії з кулькою на схилі\footnote{У цьому та наступних розділах
штрих у виразах типу $V'(\phi)$ завжди означає похідну за
скалярним полем $\phi$, тобто $V' \equiv dV/d\phi$. Це відрізняється
від §\,\ref{sec:friedmann}, де штрих позначав похідну за конформним
часом~$\eta$.}.

%% --------------------------------------------------------
\section{Динаміка повільного кочення (slow-roll) та де-Сіт\-те\-рів\-ський розв'язок}
\label{subsec:slowroll_dynamics}
%% --------------------------------------------------------

Варіюючи дію~\eqref{eq:inflaton_action} за метрикою, отримуємо тензор енергії-імпульсу скалярного поля. Для однорідного інфлатона $\phi(t)$ він набуває вигляду тензора ідеальної рідини з густиною енергії $\bar{\rho}_\phi$ та тиском $\bar{p}_\phi$:
\begin{equation}\label{eq:inflaton_rho_p}
  \bar{\rho}_\phi = \frac{\dot{\phi}^2}{2} + V(\phi),
  \qquad
  \bar{p}_\phi   = \frac{\dot{\phi}^2}{2} - V(\phi).
\end{equation}

Щоб зрозуміти, за яких умов така система здатна забезпечити інфляцію, звернемося до її фундаментального визначення: інфляція --- це епоха прискореного розширення Всесвіту, тобто $\ddot{a} > 0$. Знайдемо зв'язок цієї умови зі зміною параметра Хаббла $H \equiv \dot{a}/a$. Пряме диференціювання за часом дає:
\begin{equation*}
  \dot{H} = \frac{\ddot{a}}{a} - \left(\frac{\dot{a}}{a}\right)^2 = \frac{\ddot{a}}{a} - H^2,
  \quad \Leftrightarrow \quad
  \frac{\ddot{a}}{a} = H^2 \left( 1 - \left[-\frac{\dot{H}}{H^2}\right] \right) =  H^2 \left( 1 - \varepsilon \right) .
\end{equation*}
Звідси видно, що умова прискорення $\ddot{a} > 0$ виконується тоді й лише тоді, коли безрозмірна величина:

\begin{equation}\label{eq:eps}
    \varepsilon \equiv -\frac{\dot{H}}{H^2}
\end{equation}
є меншою за одиницю.

Цей параметр $\varepsilon$ має глибокий фізичний зміст: він характеризує дробову зміну параметра Хаббла за один $e$-fold розширення ($dN = H dt$), тобто $\varepsilon = -d\ln H / dN$. Введемо також другий безрозмірний параметр $\eta_{\mathrm{H}}$, який вимірює відносну швидкість зміни (прискорення) самого скалярного поля:
\begin{equation}\label{eq:eta}
  \eta_{\mathrm{H}} \equiv -\frac{\ddot{\phi}}{H\dot{\phi}}.
\end{equation}

Зв'яжемо кінематичні характеристики \eqref{eq:eps} та \eqref{eq:eta} з енергетичними компонентами поля~\eqref{eq:inflaton_rho_p}. Запишемо перше рівняння Фрідмана~\eqref{eq:F1}, використовуючи редуковану масу Планка $M_\mathrm{Pl} \equiv (8\pi G)^{-1/2}$, а також рівняння неперервності ($\dot{\bar{\rho}}_\phi + 3H(\bar{\rho}_\phi + \bar{p}_\phi) = 0$):
\begin{align}
  H^2 &= \frac{1}{3M_\mathrm{Pl}^2} \left(\frac{\dot\phi^2}{2} + V(\phi)\right), \label{eq:friedmann_inflaton} \\
  \dot{H} &= -\frac{1}{2M_\mathrm{Pl}^2}(\bar{\rho}_\phi + \bar{p}_\phi) = -\frac{\dot{\phi}^2}{2M_\mathrm{Pl}^2}. \label{eq:hdot_inflaton}
\end{align}

Підставляючи вираз для $\dot{H}$~\eqref{eq:hdot_inflaton} в означення параметра $\varepsilon$, ми отримуємо його кінематичний вираз через динамічні характеристики поля. Щоб наочно розкрити його фізичний зміст, виразимо цей параметр явно через енергетичні компоненти інфлатона --- кінетичну густину енергії $K_\phi \equiv \dot{\phi}^2/2$ та повну густину енергії $\bar{\rho}_\phi = 3H^2 M_\mathrm{Pl}^2$:
\begin{equation}\label{eq:epsilon_energy_ratio}
  \varepsilon = \frac{\dot\phi^2}{2H^2 M_\mathrm{Pl}^2} = 3 \cdot \frac{\tfrac{1}{2}\dot{\phi}^2}{3H^2 M_\mathrm{Pl}^2} = 3 \, \frac{K_\phi}{\bar{\rho}_\phi}.
\end{equation}
Тепер логіка стає абсолютно прозорою: оскільки за визначенням $\bar{\rho}_\phi = K_\phi + V(\phi)$, вимога повільного кочення $\varepsilon \ll 1$ (що гарантує тривалу інфляцію) фізично означає, що кінетична енергія поля становить лише малу частку від повної густини енергії, тобто є нехтовно малою порівняно з потенціалом: $K_\phi \ll V(\phi)$.

У цьому режимі, який називають \emph{наближенням повільного кочення} (slow-roll), динаміка системи радикально спрощується:
\begin{equation}\label{eq:sr_approximations}
  H^2 \approx \frac{V(\phi)}{3M_\mathrm{Pl}^2}, \qquad
  w_\phi = \frac{\bar{p}_\phi}{\bar{\rho}_\phi} = \frac{\dot\phi^2/2 - V}{\dot\phi^2/2 + V} \approx -1.
\end{equation}
Оскільки кінетична енергія мала, густина енергії залишається практично сталою протягом багатьох актів розширення: $\bar{\rho}_\phi \approx V(\phi) \approx \mathrm{const}$. Тоді з першого рівняння Фрідмана випливає, що параметр Хаббла також є квазісталим ($H \approx \mathrm{const}$), а масштабний фактор демонструє експоненційне (де-Сіттерівське) зростання:
\begin{equation}\label{eq:a_exponential}
  a(t) = a_0\,e^{Ht}.
\end{equation}
Саме цей режим розв'язує класичні космологічні проблеми: за $N \equiv \int H dt \sim 50$--$60$ $e$-folds будь-яка початкова просторова кривина розгладжується експоненціально, а причинно пов'язана область розростається до масштабів, що значно перевищують наш сучасний спостережуваний горизонт.

Для підтримання такого режиму поле має котитися плавно (без суттєвого прискорення), тому, з умови $|\eta_\mathrm{H}| \ll 1$ випливає, що сили інерції у рівнянні Кляйна-Ґордона-Фока є малими порівняно з гравітаційним тертям ($|\ddot\phi| \ll 3H|\dot\phi|$). Рівняння руху поля~\eqref{eq:klein_gordon} втрачає другий рядок і зводиться до системи першого порядку (динамічного атрактора):
\begin{equation}\label{eq:sr_attractor}
  3H\dot{\phi} \approx -V'(\phi).
\end{equation}

%Завдяки атрактору~\eqref{eq:sr_attractor}, ми можемо виразити параметри slow-roll виключно через локальну геометричну форму самого потенціалу $V(\phi)$ та його похідні:
%\begin{equation}\label{eq:slow_roll_pot}
%  \varepsilon \approx \epsilon_V
%  \equiv \frac{M_\mathrm{Pl}^2}{2}\left(\frac{V'}{V}\right)^2,
%  \qquad
%  \eta_{\mathrm{H}} \approx \eta_V
%  \equiv M_\mathrm{Pl}^2\,\frac{V''}{V}.
%\end{equation}

Завдяки атрактору~\eqref{eq:sr_attractor}, ми можемо виразити параметри slow-roll виключно через локальну геометричну форму самого потенціалу $V(\phi)$ та його похідні. Зручно ввести безрозмірні \emph{потенціальні параметри повільного кочення}:
\begin{equation}\label{eq:slow_roll_pot}
  \epsilon_V \equiv \frac{M_\mathrm{Pl}^2}{2}\left(\frac{V'}{V}\right)^2,
  \qquad
  \eta_V \equiv M_\mathrm{Pl}^2\,\frac{V''}{V}.
\end{equation}
У наближенні повільного кочення вони пов'язані з кінематичними параметрами $\varepsilon$ та $\eta_{\mathrm{H}}$ (означеними в~\eqref{eq:eps}--\eqref{eq:eta}) як:
\begin{equation}\label{eq:eps_eta_V}
  \varepsilon \approx \epsilon_V,
  \qquad
  \eta_{\mathrm{H}} \approx \epsilon_V - \frac{1}{2}\eta_V.
\end{equation}

Доведемо тотожність для $\epsilon_V$, підставивши швидкість на атракторі~\eqref{eq:sr_attractor} у вираз~\eqref{eq:epsilon_energy_ratio}:
\begin{equation*}
  \varepsilon = \frac{\dot\phi^2}{2H^2 M_\mathrm{Pl}^2} \approx \frac{(-V'/3H)^2}{2H^2 M_\mathrm{Pl}^2} = \frac{(V')^2}{18 H^4 M_\mathrm{Pl}^2}.
\end{equation*}
Враховуючи slow-roll значення параметра Хаббла $H^2 \approx V / (3M_\mathrm{Pl}^2)$, остаточно маємо:
\begin{equation*}
  \varepsilon \approx \frac{(V')^2}{18 M_\mathrm{Pl}^2} \left( \frac{3M_\mathrm{Pl}^2}{V} \right)^2 = \frac{M_\mathrm{Pl}^2}{2}\left(\frac{V'}{V}\right)^2 \equiv \epsilon_V.
\end{equation*}

Таким чином, для існування інфляції потенціал скалярного поля має бути достатньо пологим ($\epsilon_V < 1$, що забезпечує прискорене розширення $\ddot{a} > 0$), а для її тривалості --- надзвичайно пологим ($\epsilon_V \ll 1$ та $|\eta_V| \ll 1$).

Інфляція невідворотно завершується, коли поле виходить на крутий схил потенціалу, де виконується умова $\varepsilon(\phi_{\mathrm{end}}) = 1$ (поширена помилка --- вважати, що це відбувається в мінімумі; насправді $\varepsilon = 1$ відповідає рівності $K_\phi = \frac{1}{2}V$, що реалізується ще на схилі).

У цій точці прискорене розширення простору миттєво припиняється ($\ddot{a}=0$), проте саме скалярне поле, набувши колосальної швидкості, не зупиняється: за інерцією воно мчить далі вниз по схилу, проскакує мінімум потенціалу $V(\phi)$ і піднімається по протилежному схилу, де на мить завмирає й починає рух у зворотному напрямку.

Внаслідок дії хаббловського тертя $3H\dot{\phi}$, яке відіграє роль в'язкого опору середовища, амплітуда цього руху стрімко згасає, і поле переходить у режим високочастотних когерентних коливань навколо мінімуму.

Кількісним критерієм «високочастотності» на цьому етапі є виконання строгої нерівності:
\begin{equation*}
    \omega \approx m_\phi \gg H,
\end{equation*}
де $\omega = \sqrt{V''(\phi_{\mathrm{min}})} \equiv m_\phi$ --- власна частота коливань поля (яка визначається кривиною потенціалу в мінімумі й тотожна масі квантів інфлатона), а $H$ --- параметр Хаббла. Оскільки хабблівський час $\tau_H \sim H^{-1}$ задає характерний масштаб розширення простору, виконання умови $m_\phi \gg H$ означає, що за один такий такт Всесвіту поле встигає здійснити тисячі або мільйони повних коливань ($\tau_{\mathrm{osc}} \ll \tau_H$). На тлі динаміки метрики ці осциляції є надзвичайно швидкими, що дозволяє усереднювати рівняння стану поля за період коливань. Саме цей коливальний етап запускає процес генерації матерії, який ми детально розглянемо у розділі~\ref{sec:reheating}.



%% --------------------------------------------------------
\section{Хаотична інфляція Лінде}
\label{sec:chaotic}
%% --------------------------------------------------------

Лінде~\cite{linde_1983} запропонував радикально простий погляд
на початкові умови: значення $\phi$ у різних точках раннього
Всесвіту є хаотичним~--- немає потреби в особливому «стартовому»
стані. Найпростіший потенціал:
\begin{equation*}
  V(\phi) = \tfrac{1}{2}m^2\phi^2.
\end{equation*}

Де $\phi$ велике і $\varepsilon \ll 1$~--- там інфляція іде довго
і виникає величезна причинно зв'язана «бульбашка». Де $\phi$ мале~---
інфляція швидко завершується. Наш спостережуваний Всесвіт~---
одна з таких бульбашок.

\textbf{Структура всередині однієї бульбашки.} Коли інфляція
в нашій бульбашці завершилась, простір всередині неї вже був
практично однорідним~--- розширення загладило всі початкові
неоднорідності $\phi$. Але «практично» не означає «абсолютно»:
залишились крихітні квантові флуктуації, заморожені під час
інфляції. Саме вони стали насінням галактик і є тим, що ми
бачимо на карті CMB.

Всередині однієї бульбашки існує єдиний причинно зв'язаний
простір-час~--- один Всесвіт. Але його розмір після $N \sim 60$
$e$-folds становить $\sim e^{60}$ разів більше за наш
спостережуваний горизонт ($\sim 46$~Гпк). Тому навіть у
межах «нашої бульбашки» існують області, з якими ми ніколи
не матимемо причинного контакту~--- вони назавжди за
нашим горизонтом подій.

\textbf{Між бульбашками.} Простір між різними бульбашками
продовжує розширюватись інфляційно~--- швидше за світло.
Дві бульбашки, що утворились незалежно, розділені областю
вічної інфляції і принципово не можуть обмінятись сигналом.
Це не просто «дуже далеко»~--- між ними немає спільного
майбутнього світлового конуса. У строгому сенсі загальної теорії відносності вони
належать різним причинно відокремленим регіонам
одного глобального простору-часу, а не різним «всесвітам»~---
хоча в популярній літературі їх так і називають.

\textbf{Вічна інфляція.} Квантові флуктуації інфлатона на
суперхабблових масштабах можуть локально \emph{збільшити} $\phi$,
повернувши його назад на схил, якщо флуктуація достатньо велика.
Класичний зсув поля за один $e$-fold: $|\delta\phi_{\rm class}|
\sim |\dot\phi|/H \approx \sqrt{2\varepsilon}\,M_{\rm Pl}$.
Квантова флуктуація за той самий час: $\delta\phi_{\rm quant}
\sim H/(2\pi)$. Умова «самовідтворення»~---
квантова флуктуація перевищує класичне зміщення:
\begin{equation*}
  \frac{H}{2\pi} \gtrsim \sqrt{2\varepsilon}\,M_{\rm Pl}.
\end{equation*}
Для типових моделей $\varepsilon \lesssim 1$, тому достатньою
умовою є $H^2 \gtrsim M_{\rm Pl}^2$, що за допомогою
$H^2 \approx V/(3M_{\rm Pl}^2)$ дає
\begin{equation*}
  V(\phi) \gtrsim 3M_{\rm Pl}^4.
\end{equation*}
У таких областях інфляція ніколи не зупиняється глобально:
поки одна бульбашка завершує інфляцію і термалізується,
сусідні регіони продовжують розширюватись, породжуючи нові
бульбашки. Це \textbf{вічна інфляція}~--- мультивсесвіт
як математичний наслідок, а не постулат.

Строгий теоретичний опис цього динамічного режиму забезпечується стохастичним
підходом Старобінського~\cite{starobinsky_1986}. У цій формалізації квантові флуктуації короткохвильових
мод розглядаються як випадкове джерело (білий шум), що діє на класичну
довгохвильову координату поля. Еволюція інфлатона набуває вигляду стохастичного
рівняння Ланжевена, де класичний дрейф по схилу потенціалу конкурує з квантовою
дифузією. Коли поле перебуває поблизу планківської межі, дифузійний внесок
повністю домінує, що призводить до квантового блукання поля вгору по потенціалу
і робить процес генерації нових областей простору просторово нескінченним.

\textbf{Статус.} Потенціал $\tfrac{1}{2}m^2\phi^2$ передбачає
$r \approx 8/N \approx 0.13$~--- вище поточної межі BICEP/Keck
($r < 0.036$, див. додаток~\ref{app:data}), тому ця конкретна модель виключена.
Але ідея хаотичних початкових умов і вічної інфляції
залишається живою у більш пологих потенціалах.
Сучасним переосмисленням парадигми Лінде є клас моделей
$\alpha$-аттракторів~\cite{ferrara_2013,kallosh_2013}
та інфляція з немінімальним зв'язком інфлатона з кривиною простору-часу.
У цих сценаріях при великих значеннях поля потенціал природно виходить на
асимптотичне плато, що дозволяє зберегти концепцію хаотичного старту високої
енергії, але радикально знижує амплітуду тензорних збуджень $r$, повністю
узгоджуючи теорію з останніми даними Planck та BICEP/Keck.

Разом з тим жодна з «успішних» моделей~--- тих, що проходять
обмеження Planck на $n_s$ і $r$~--- не передбачає нічого
особливого на малих масштабах $k \gtrsim 1\,h/$Мпк. Саме там
первинний спектр залишається практично необмеженим
спостереженнями. І саме там у 2022--2023 роках телескоп JWST
зафіксував надлишок масивних галактик при $z \sim 10$--$12$:
об'єктів, які за стандартним плоским спектром просто не встигли
б вирости за відведений час~\cite{finkelstein_2023,boylan_kolchin_2023}.

Природне пояснення~--- локальне підсилення первинного спектра
саме на цих масштабах. Якщо потенціал інфлатона має особливість
(сходинку або ділянку ультра-повільного кочення) у потрібному
місці, флуктуації, що виходять за горизонт в цей момент,
зазнають різкого підсилення~--- формується \emph{зламаний спектр
потужності} (рис.~\ref{fig:jwst_feature}). Чи здатна ця ідея витримати зіставлення з повним
масивом даних Planck і JWST~--- саме це є предметом
наступного розділу.

%% --------------------------------------------------------
%% РИСУНОК: МОДИФІКАЦІЯ ПОТЕНЦІАЛУ ТА ЗЛАМАНИЙ СПЕКТР
%% --------------------------------------------------------
\begin{figure}[htbp]
  \centering
  \begin{tikzpicture}
    \pgfplotsset{compat=1.18}

    % Лівий графік: Потенціал з особливістю (USR)
    \begin{axis}[
        width=0.48\textwidth,
        height=5.5cm,
        axis lines=left,
        xlabel={$\phi$},
        ylabel={$V(\phi)$},
        ymin=0, ymax=4.0,
        xmin=0, xmax=4.5,
        ticks=none,
        title={Потенціал інфлатона}
      ]
      % Монотонна функція з вираженим плато
      \addplot[thick, color=blue, domain=0:4.2, samples=100, smooth]
        { (2.5 / (1 + exp(-3*(x-2)))) + 0.3*x };

      % Позначення ділянки особливості (USR)
      \draw[dashed, gray] (axis cs:1.3,0) -- (axis cs:1.3,1.5);
      \draw[dashed, gray] (axis cs:2.7,0) -- (axis cs:2.7,1.5);
      \draw[<->, >=stealth, gray] (axis cs:1.3,0.4) -- node[above, font=\scriptsize, text=black] {USR} (axis cs:2.7,0.4);

      % Акуратна коротка стрілка-вказівка
      \node[above, font=\scriptsize] at (axis cs:1.6,2.7) {Особливість (сходинка)};
      \draw[->, >=stealth, thick, darkgray] (axis cs:1.6,2.6) -- (axis cs:2.0,1.55);
    \end{axis}

    % Правий графік: Зламаний спектр потужності
    \begin{axis}[
        width=0.48\textwidth,
        height=5.5cm,
        xshift=6.8cm,
        axis lines=left,
        xlabel={$\ln k$},
        ylabel={$\ln \mathcal{P}_\zeta(k)$},
        ymin=-1, ymax=4.5,
        xmin=0, xmax=5,
        ticks=none,
        title={Первинний спектр потужності}
      ]
      % Масштаби CMB (плато) та різкий пік на малих масштабах
      \addplot[thick, color=red, domain=0:5, samples=150, smooth]
        {0.5 + 3.2*exp(-(x-3.4)^2/0.18)};

      % Позначення зон обмежень
      \draw[dashed, gray] (axis cs:0,0.5) -- (axis cs:2.2,0.5);
      \node[below, font=\scriptsize, text=darkgray] at (axis cs:1.1,0.5) {CMB (Planck)};

      \draw[dashed, gray] (axis cs:3.4,-1) -- (axis cs:3.4,3.7);
      \node[above, font=\scriptsize, text=red] at (axis cs:3.4,3.7) {Пік JWST};
    \end{axis}
  \end{tikzpicture}
  \caption{Схематичне зображення механізму генерації зламаного спектра.
           Ліворуч: потенціал інфлатона з короткою ділянкою ультра-повільного кочення
           (USR --- ultra-slow roll), де класична швидкість поля падає.
           Праворуч: відповідний спектр потужності кривини $\mathcal{P}_\zeta(k)$,
           що зберігає масштабну інваріантність на масштабах CMB,
           але має локальний надлишок амплітуди на високих хвильових числах.}
  \label{fig:jwst_feature}
\end{figure}


%% --------------------------------------------------------
\section{Форма потенціалу та обмеження спостережень}
%% --------------------------------------------------------

Спостереження CMB дають нам інформацію не про весь потенціал
$V(\phi)$, а лише про надзвичайно вузький його проміжок. Масштаби,
зафіксовані на картах анізотропії Planck, виходили за горизонт
протягом $\sim\!7 - 8$ $e$-folds із загальних $50 - 60$. Іншими
словами, кожне кутове масштабне відхилення температури CMB
відповідає певному значенню $\phi$ під час інфляції, і вся ця
«пам'ять» закодована лише у вузькій ділянці потенціалу~---
тій, що залишила слід на кутовому спектрі анізотропій
(див.\ рис.~\ref{fig:cmb_tt}).

%\begin{figure}[htbp]
%\centering
%\includegraphics[width=\linewidth]{planck_cmb_map.png}
%\caption{Карта анізотропій CMB від місії Planck~2018
%         \cite{planck_2018}. Кольори відображають відхилення
%         температури від середнього значення $T_0 = 2.725\,\mathrm{K}$
%         в діапазоні $\pm 300\,\mu\mathrm{K}$. Саме ця карта
%         є спостережуваним відбитком первинних флуктуацій
%         інфлатона.}
%\label{fig:planck_map}
%\end{figure}

Сьогодні існує понад 100 конкурентних моделей інфляції.
Більшість із них дають практично однаковий фіт для великих
кутових масштабів, де Planck найточніший. Різниця між
моделями виявляється на малих масштабах
($k \gg 1\,\mathrm{Mpc}^{-1}$), де поточні спостереження
мають значно слабші обмеження.

Таблиця~\ref{tab:inflation_models} порівнює передбачення
найпопулярніших класів моделей для двох вимірюваних
параметрів~--- $n_s$ і $r$~--- з поточними обмеженнями
Planck~2018 \cite{planck_2018} + BICEP/Keck \cite{bicep_keck_2021}.

\begin{table}[htbp]
\centering
\caption{Передбачення основних класів інфляційних моделей.
         Значення наведені для $N = 50$--$60$ $e$-folds до
         кінця інфляції. Спостережувані межі:
         $n_s = 0.9649 \pm 0.0042$,\ $r < 0.036$
         \cite{planck_2018, bicep_keck_2021}.}
\label{tab:inflation_models}
\begin{tblr}{
  colspec = {XXXXX[l,m]},
  row{1}  = {c, m, font=\bfseries},
}
\toprule
Модель & Потенціал $V(\phi)$ & $n_s$ & $r$ & Статус \\
\midrule
Старобінський ($R^2$) \cite{starobinsky_1980}
  & плато & $1 - \frac{2}{N}$ & $\frac{12}{N^2}$
  & узгоджується (детально~--- §\,\ref{sec:starobinsky}) \\
Хіггсова інфляція \cite{bezrukov_2008}
  & плато (нем.~зв'язок) & $\approx 0.965$ & $\approx 0.004$
  & узгоджується \\
Natural inflation \cite{freese_1990}
  & $1 - \cos(\phi/f)$ & $\approx 0.95$--$0.97$ & $0.01$--$0.1$
  & під питанням \\
Хаотична, $\phi^2$ \cite{linde_1983}
  & $\frac{1}{2}m^2\phi^2$ & $1 - \frac{2}{N}$ & $\frac{8}{N}$
  & виключена ($r \sim 0.13$) \\
Хаотична, $\phi^{4/3}$ \cite{linde_1983}
  & $\lambda\phi^{4/3}$ & $1 - \frac{5}{3N}$ & $\frac{16}{3N}$
  & під питанням \\
Зі сходинкою
  & плато + $\delta V$ & залежить від $k_{\rm break}$ & мале
  & перевіряється \\
\bottomrule
\end{tblr}
\end{table}

Серед усіх моделей таблиці~\ref{tab:inflation_models} модель Старобінського
займає особливе місце~--- вона не просто «узгоджується», а потрапляє в самий
центр довірчої області Planck. Більш того, вона побудована без введення нового
поля «руками»~--- лише через модифікацію геометричної частини дії. Розберемо
цей механізм детально.

%% --------------------------------------------------------
\section{Модель Старобінського та $f(R)$ гравітація}
\label{sec:starobinsky}
%% --------------------------------------------------------

Історично перша модель інфляції, запропонована О.~Старобінським у
1980~році~\cite{starobinsky_1980}, базується не на введенні нового
скалярного поля «руками», а на модифікації самої геометричної частини
дії Айнштайна--Гільберта.

%% --------------------------------------------------------
\subsection*{Що таке $f(R)$ гравітація?}
%% --------------------------------------------------------

\begin{keybox}{Пропедевтика}
Стандартна ЗТВ будується на принципі
найменшої дії з дією Айнштайна--Гільберта:
\[
  S_\mathrm{EH} = \frac{M_\mathrm{Pl}^2}{2}\int d^4x\sqrt{-g}\,R,
\]
де $R$~--- скалярна кривизна Річчі. Ця дія \emph{лінійна} по $R$.
Природне узагальнення~--- замінити $R$ на довільну функцію $f(R)$:
\[
  S = \frac{M_\mathrm{Pl}^2}{2}\int d^4x\sqrt{-g}\,f(R).
\]
\end{keybox}

Запитання: чи це справді «нова» теорія гравітації,
чи просто ЗТВ із прихованим скалярним полем?

Відповідь дає фундаментальна теорема про конформну еквівалентність.

%% --------------------------------------------------------
\subsection*{Теорема: $f(R) \;\leftrightarrow\; $ ЗТВ + скалярне поле}
%% --------------------------------------------------------

\begin{keybox}{Теорема (конформна еквівалентність)}
Будь-яка $f(R)$-теорія гравітації з $f''(R)\neq 0$ у
\emph{йорданівському фреймі} (Jordan frame) математично еквівалентна
стандартній ЗТВ з масивним скалярним полем $\phi$ у
\emph{айнштайнівському фреймі} (Einstein frame).
\end{keybox}

\textbf{Виведення (покрокове).}

\medskip
\noindent\textit{Крок 1. Введення допоміжного поля.}
Запишемо $f(R)$-дію через допоміжне поле $\chi$:
\begin{equation}\label{eq:fR_auxiliary}
  S = \frac{M_\mathrm{Pl}^2}{2}\int d^4x\sqrt{-g}
      \bigl[f(\chi) + f'(\chi)(R - \chi)\bigr].
\end{equation}
Рівняння руху для $\chi$: варіювання за $\chi$ дає
$f''(\chi)(R-\chi) = 0$.
За умови $f''(\chi)\neq 0$ це означає $\chi = R$~---
тобто дія~\eqref{eq:fR_auxiliary} точно відтворює вихідну $f(R)$-теорію.

\medskip
\noindent\textit{Крок 2. Перейменування та структура.}
Введемо скалярне поле $\varphi \equiv f'(\chi) = f'(R)$.
Тоді дію~\eqref{eq:fR_auxiliary} можна переписати як:
\begin{equation}\label{eq:fR_jordan}
  S = \frac{M_\mathrm{Pl}^2}{2}\int d^4x\sqrt{-g}
      \bigl[\varphi R - U(\varphi)\bigr],
  \qquad
  U(\varphi) \equiv \chi(\varphi)\,\varphi - f(\chi(\varphi)).
\end{equation}
Це скалярно-тензорна теорія Бранса--Дікке з $\omega_\mathrm{BD}=0$
у йорданівському фреймі: поле $\varphi$ нелінійно зачеплене з
кривизною через доданок $\varphi R$.

\begin{keybox}{Що таке фрейми?}
\emph{Йорданівський фрейм}~--- «фізичний» фрейм, у якому
частинки матерії рухаються геодезичними метрики $g_{\mu\nu}$,
а скалярне поле зачеплене з кривизною.
\emph{Айнштайнівський фрейм}~--- фрейм, у якому гравітаційна
частина дії має стандартний вигляд Айнштайна--Гільберта з $M_\mathrm{Pl}^2 R/2$,
але скалярне поле з'являється явно з кінетичним членом.
Перехід між фреймами~--- це конформне перетворення метрики
$\tilde g_{\mu\nu} = \Omega^2(\varphi)\,g_{\mu\nu}$,
яке є \emph{математичною} переформулювання, а не фізичним перетворенням.
\end{keybox}

\medskip
\noindent\textit{Крок 3. Конформне перетворення до айнштайнівського фрейму.}
Виконаємо конформне перетворення метрики:
\begin{equation}\label{eq:conformal_transform}
  \tilde{g}_{\mu\nu} = \varphi\, g_{\mu\nu},
  \qquad
  \sqrt{-\tilde{g}} = \varphi^2\sqrt{-g},
  \qquad
  \tilde{R} = \frac{1}{\varphi}\left(R - \frac{3\Box\varphi}{\varphi}
  + \frac{3(\partial\varphi)^2}{2\varphi^2}\right).
\end{equation}
Після підстановки та інтегрування частинами дія~\eqref{eq:fR_jordan}
набуває канонічного вигляду в айнштайнівському фреймі:
\begin{equation}\label{eq:fR_einstein}
  S = \int d^4x\sqrt{-\tilde{g}}\left[
      \frac{M_\mathrm{Pl}^2}{2}\tilde{R}
      - \frac{1}{2}(\tilde\nabla\phi)^2
      - V(\phi)
      \right],
\end{equation}
де канонічне скалярне поле $\phi$ і його потенціал $V(\phi)$
визначаються через:
\begin{equation}\label{eq:fR_canonical_field}
  \phi \equiv \sqrt{\frac{3}{2}}\,M_\mathrm{Pl}\ln\varphi
  = \sqrt{\frac{3}{2}}\,M_\mathrm{Pl}\ln f'(R),
  \qquad
  V(\phi) = \frac{M_\mathrm{Pl}^2}{2}
            \frac{\chi f'(\chi) - f(\chi)}{[f'(\chi)]^2}.
\end{equation}

\begin{keybox}{Фізичний зміст}
Доданок $R^2$ у дії~--- це не просто «виправлення»:
він є \emph{кінетичним членом для нового ступеня вільності}.
$f(R)$-гравітація має на один ступінь вільності більше ніж ЗТВ
саме через те, що $f''(R)\neq 0$: четвертий порядок рівнянь поля
розщеплюється на другий порядок для метрики і другий порядок
для скалярона $\phi$.
При $f''(R) = 0$ (тобто $f(R) = R$)~--- звичайна ЗТВ,
новий ступінь вільності зникає.
\end{keybox}

%% --------------------------------------------------------
\subsection*{Застосування до моделі Старобінського}
%% --------------------------------------------------------

\textbf{Вихідна дія.} Модель Старобінського:
\begin{equation}\label{eq:starobinsky_action}
  f(R) = R + \frac{R^2}{6M^2},
\end{equation}
де $M$~--- характерний масштаб, що фіксується амплітудою первинного
спектра $A_s$.

\textbf{Потенціал скалярона.}
Застосовуємо формулу~\eqref{eq:fR_canonical_field}:
$f'(R) = 1 + R/(3M^2)$, тому
$\varphi = 1 + R/(3M^2)$ і $R = 3M^2(\varphi-1)$.
Підставляючи у~\eqref{eq:fR_canonical_field}:
\begin{equation}\label{eq:starobinsky_potential}
  V(\phi) = \frac{3M^2 M_\mathrm{Pl}^2}{4}
            \left(1 - e^{-\sqrt{2/3}\,\phi/M_\mathrm{Pl}}\right)^2.
\end{equation}

%---------------------------------------------------------
\begin{figure}[h!]
\centering
\begin{tikzpicture}
  \begin{axis}[
    width=0.7\textwidth, height=5cm,
    axis lines=left,
    xlabel={$\phi/M_{\mathrm{Pl}}$}, ylabel={$V(\phi)/V_0$},
    xmin=0, xmax=5, ymin=0, ymax=1.2,
    ticks=none,
    title={Потенціал скалярона (модель Старобінського)}
  ]
    \addplot[ultra thick, color=accent, domain=0:5, samples=100]
      {(1 - exp(-sqrt(2/3)*x))^2};
    \draw[dashed, gray] (0,1) -- (5,1)
      node[right, font=\tiny] {$V_0 = \tfrac{3}{4}M^2M_\mathrm{Pl}^2$};
    \node[anchor=south east, font=\scriptsize, text=gray]
      at (axis cs:4.2,0.85) {slow-roll область};
    \draw[->, gray, thin] (axis cs:3.8,0.82) -- (axis cs:2.5,0.7);
  \end{axis}
\end{tikzpicture}
\caption{Потенціал скалярона~\eqref{eq:starobinsky_potential}.
На великих $\phi$ потенціал виходить на плато $V\to V_0$~---
саме там відбувається slow-roll інфляція. На відміну від хаотичної
інфляції ($V\propto\phi^2$), плато забезпечує дуже мале тензорне
відношення $r$.}
\end{figure}
%---------------------------------------------------------

\textbf{Slow-roll параметри і передбачення.}
На плато ($\phi \gg M_\mathrm{Pl}$) параметри slow-roll~\eqref{eq:slow_roll_potential}:
\begin{align}
  \varepsilon_V &\approx \frac{4}{3}\,e^{-2\sqrt{2/3}\,\phi/M_\mathrm{Pl}}
                \approx \frac{3}{4N^2}, \\
  \eta_V        &\approx -\frac{4}{3}\,e^{-\sqrt{2/3}\,\phi/M_\mathrm{Pl}}
                \approx -\frac{1}{N}.
\end{align}
Це дає спостережувані параметри (для $N = 50$--$60$):
\begin{equation}
  \boxed{n_s \approx 1 - \frac{2}{N} \approx 0.967,}
  \qquad
  \boxed{r \approx \frac{12}{N^2} \approx 0.003.}
\end{equation}
Обидва значення потрапляють у самий центр довірчої області
Planck~2018~\cite{planck_2018} і BICEP/Keck~\cite{bicep_keck_2021}~---
саме тому модель Старобінського є сьогодні «золотим стандартом» інфляції.

\medskip
\textbf{Масштаб $M$.}
Амплітуда первинного спектра $A_s \approx 2.1\times10^{-9}$ фіксує:
\begin{equation}
  M \approx 1.3\times10^{-5}\,M_\mathrm{Pl} \approx 3\times10^{13}\,\text{ГеВ}.
\end{equation}
Це природний масштаб теорій Великого об'єднання (GUT)~---
ще одна підказка що інфляція і GUT-фізика пов'язані.

\begin{keybox}{Важливе концептуальне зауваження: який фрейм фізичний?}
Теорема конформної еквівалентності говорить що $f(R)$ і
ЗТВ+скалярне поле \emph{математично} еквівалентні.
Але чи еквівалентні вони \emph{фізично}~---
досі відкрите концептуальне питання.
У йорданівському фреймі маси частинок~--- константи,
але гравітаційна «константа» $G_\mathrm{eff}\propto 1/\varphi$ змінна.
В айнштайнівському фреймі $G$ стала, але маси частинок залежать від $\phi$.
Якщо матерія мінімально зачеплена з метрикою (що є стандартним
припущенням), обидва фрейми дають однакові спостережувані
на рівні класичної теорії~--- але можуть розрізнятися
при квантових поправках.
Для розрахунку спектра $P_{\mathcal{R}}(k)$ це не важливо:
результат однаковий в обох фреймах.
\end{keybox}


Якщо геометрія потенціалу на цій «непомітній» ділянці різко
змінюється~--- наприклад, з'являється локальна сходинка або ділянка
ультра-повільного кочення (USR)~--- режим slow-roll тимчасово
порушується. Флуктуації, що виходять за горизонт саме в цей момент,
зазнають різкого підсилення, формуючи \emph{зламаний спектр
потужності}. Цей механізм відкриває широке поле для феноменологічного
пошуку: окрім вирішення проблеми надлишку ранніх структур за даними JWST,
тут є над чим попрацювати в контексті напруженості Хаббла (Hubble tension).
Оскільки модифікація первинного спектра на малих масштабах неминуче
зміщує найкращі фіти фонових параметрів $\Lambda$CDM при сумісному MCMC-аналізі,
можливо існує реальна перспектива часткового послаблення цієї кризи, що
детально розглядається у §\,\ref{sec:future_research}. Але перш ніж перейти
до нього~--- нам потрібно зрозуміти, як саме квантові флуктуації
інфлатона перетворюються на початкові неоднорідності Всесвіту.
Саме тут інфляція є унікальною: на відміну від будь-якого
сценарію з довільними початковими умовами, вона генерує ці
неоднорідності з квантового вакууму~--- без жодного зовнішнього
задання.
Як саме це відбувається --- предмет наступного розділу.



% ============================================================
% Задачі до розділу 2: Інфлатон і його потенціал
\section*{Задачі}
\addcontentsline{toc}{section}{Задачі}
% ============================================================

\begin{problem}{Slow-roll параметри та умова інфляції}
Для скалярного поля $\phi$ із потенціалом $V(\phi)$ у наближенні повільного кочення:

\begin{enumerate}
    \item Використовуючи $\varepsilon = -\dot{H}/H^2$ та рівняння Фрідмана,
          доведіть, що $\varepsilon = \dfrac{\dot{\phi}^2}{2H^2 M_{\mathrm{Pl}}^2}$.
    \item Покажіть, що умова інфляції $\ddot{a} > 0$ еквівалентна $\varepsilon < 1$.
    \item Виразіть $\varepsilon$ через потенціал $V(\phi)$ та його похідну $V'(\phi)$
          у наближенні slow-roll, тобто виведіть $\epsilon_V$.
\end{enumerate}
\end{problem}

% -------------------------------------------------------------------

\begin{problem}{Рівняння атрактора та його фізичний зміст}
Рівняння Кляйна--Ґордона для однорідного поля в розширюваному Всесвіті:
\[
\ddot{\phi} + 3H\dot{\phi} + V'(\phi) = 0.
\]

\begin{enumerate}
    \item У наближенні повільного кочення нехтують членом $\ddot{\phi}$.
          Чому це виправдано за умови $|\eta_H| \ll 1$?
    \item Запишіть рівняння атрактора $3H\dot{\phi} \approx -V'(\phi)$.
          Поясніть фізичну аналогію з \enquote{кулькою на схилі у в'язкому середовищі}.
    \item Використовуючи атрактор, покажіть, що $\dot{\phi}$ визначається локальним
          нахилом потенціалу, а не початковими умовами. Чому це важливо для
          передбачуваності інфляції?
\end{enumerate}
\end{problem}

% -------------------------------------------------------------------

\begin{problem}{Зв'язок між $\eta_H$ та $\eta_V$}
Кінематичний параметр $\eta_H = -\ddot{\phi}/(H\dot{\phi})$ і потенціальний параметр
$\eta_V = M_{\mathrm{Pl}}^2 V''/V$ пов'язані у наближенні slow-roll.

\begin{enumerate}
    \item Продиференціюйте рівняння атрактора $3H\dot{\phi} \approx -V'(\phi)$ за часом.
    \item Враховуючи $\dot{H} = -\dot{\phi}^2/(2M_{\mathrm{Pl}}^2)$ та $H^2 \approx V/(3M_{\mathrm{Pl}}^2)$,
          виведіть співвідношення:
          \[
          \eta_H \approx \epsilon_V - \frac{1}{2}\eta_V.
          \]
    \item Чому не можна стверджувати, що $\eta_H \approx \eta_V$? Наведіть контрприклад
          (наприклад, потенціал $V(\phi) = \frac{1}{2}m^2\phi^2$).
\end{enumerate}
\end{problem}

% -------------------------------------------------------------------

\begin{problem}{Тривалість інфляції та кількість $e$-folds}
Кількість $e$-folds інфляції визначається як $N = \int H\,dt$.

\begin{enumerate}
    \item Використовуючи $\dot{\phi} \approx -V'/(3H)$ та $H^2 \approx V/(3M_{\mathrm{Pl}}^2)$,
          покажіть, що
          \[
          dN = \frac{H}{\dot{\phi}}\,d\phi \approx \frac{1}{M_{\mathrm{Pl}}^2}\frac{V}{V'}\,d\phi.
          \]
    \item Для потенціалу $V(\phi) = \frac{1}{2}m^2\phi^2$ обчисліть $N$ як функцію
          початкового значення поля $\phi_{\text{start}}$ та кінцевого $\phi_{\text{end}}$.
    \item Якщо інфляція закінчується при $\phi_{\text{end}} \sim M_{\mathrm{Pl}}$,
          знайдіть $\phi_{\text{start}}$, необхідне для $N \approx 60$.
\end{enumerate}
\end{problem}

% -------------------------------------------------------------------

\begin{problem}{Спектральний індекс $n_s$ та його залежність від потенціалу}
Спектральний індекс скалярних збурень у наближенні slow-roll:
\[
n_s - 1 = -2\epsilon_V - \eta_V.
\]

\begin{enumerate}
    \item Для потенціалу $V(\phi) = \frac{1}{2}m^2\phi^2$ (хаотична інфляція)
          обчисліть $\epsilon_V$, $\eta_V$ та $n_s$ як функцію $N$.
    \item Чому така модель сьогодні вважається виключеною спостереженнями?
          (Використайте $n_s = 0.9649 \pm 0.0042$ та $r = 16\epsilon_V$.)
    \item Для моделі Старобінського (плато) $n_s \approx 1 - 2/N$.
          Оцініть $n_s$ для $N = 60$ та порівняйте з даними Planck.
\end{enumerate}
\end{problem}

% -------------------------------------------------------------------

\begin{problem}{Чому поле Хіггса не може бути інфлатоном?}
Потенціал Хіггса Стандартної моделі має вигляд
\[
V(h) = \frac{\lambda}{4}\left(h^2 - v^2\right)^2,
\]
де $v \approx 246$~GeV --- вакуумне очікуване значення, $\lambda \approx 0.13$,
а $M_{\mathrm{Pl}} \approx 2.4 \times 10^{18}$~GeV.

\begin{enumerate}
    \item При $h \gg v$ спростіть потенціал до вигляду $V(h) \approx \frac{\lambda}{4}h^4$.
    \item Використовуючи визначення $\epsilon_V = \dfrac{M_{\mathrm{Pl}}^2}{2}\left(\dfrac{V'}{V}\right)^2$,
          обчисліть $\epsilon_V$ як функцію $h$.
    \item Якщо інфляція відбувається при $h \sim M_{\mathrm{Pl}}$, то яке значення має $\epsilon_V$?
          Чи виконується умова інфляції $\epsilon_V < 1$? Зробіть висновок, чи може стандартне
          поле Хіггса виступати в ролі інфлатона.
    \item Відомо, що додавання до дії члена $\xi h^2 R$ (не мінімальний зв'язок із гравітацією)
          може зробити поле Хіггса придатним для інфляції. Чому цей член змінює ситуацію?
\end{enumerate}
\end{problem}

% -------------------------------------------------------------------

\begin{problem}{Хаотична інфляція Лінде (модель $V = \frac{1}{2}m^2\phi^2$)}
Для потенціалу $V(\phi) = \frac{1}{2}m^2\phi^2$:

\begin{enumerate}
    \item Обчисліть параметри повільного кочення $\epsilon_V$, $\eta_V$ та
          спектральний індекс $n_s$ як функцію $N$.
    \item Використовуючи $A_s = \dfrac{V}{24\pi^2 M_{\mathrm{Pl}}^4 \epsilon_V} \approx 2.1\times10^{-9}$
          та $N \approx 60$, оцініть масу $m$ інфлатона в одиницях $M_{\mathrm{Pl}}$ та в ГеВ.
    \item Обчисліть тензорно-скалярне відношення $r = 16\epsilon_V$ для $N = 60$.
          Порівняйте з поточною верхньою межею $r < 0.036$ (BICEP/Keck 2021).
          Чому ця модель вважається виключеною?
    \item (Бонус) Умова вічної інфляції: $V(\phi) \gtrsim 3M_{\mathrm{Pl}}^4$.
          Оцініть критичне значення $\phi$, за якого виникає вічна інфляція.
\end{enumerate}
\end{problem}